\documentclass[../main.tex]{subfiles}
\begin{document}
%==========================================TIÊU ĐỀ TẬP========================================
\begin{table}[h]
  \begin{adjustwidth}{-1cm}{}
    \begin{tabular}{>{\centering\arraybackslash}p{6cm}|>{\raggedright\arraybackslash}p{12.5cm}}
      \multirow{5}{6cm}
      % Cột bên trái: ÉPISODE và số tập
      {\\[-40pt]
        \flaregothic\fontsize{20pt}{20pt}\selectfont \centering
        {\contour{errorfour}{\textcolor{black}{ÉPISODE}}}\color{black}\\
        \burbank\fontsize{72pt}{72pt}\selectfont
        \includegraphics[height=3.12359cm]{../../../../Image/Book?/number/num21.png}
        \\[18pt]
        % \flaregothic\fontsize{14pt}{14pt}\selectfont
        % \color{epattr}BIENVENUE, \uline{\color{Banpire} Mel} !
        % \flaregothic\fontsize{18pt}{18pt}\selectfont
        % \color{epattr}ÉPISODE PERDU
        \color{black}
      }
       &
      % Dòng thứ nhất: tên tập tiếng Nhật
      {\fotskip\fontsize{18pt}{18pt}\selectfont \ruby{後}{こう}\ruby{悔}{かい}の\ruby{日}{ひ}々}
      \\[6pt]
      % Dòng thứ hai: tên tập tiếng Anh
       & {\cafeteria\fontsize{18pt}{18pt}\selectfont Days of Regret}                \\[4pt]
      % Dòng thứ ba: tên tập tiếng Việt
       & {\vnmsans\fontsize{18pt}{18pt}\selectfont Những ngày trong niềm hối hận}   \\[8pt]
      % Dòng thứ tư: Nhân vật xuất hiện
       & {\caviar{\fontsize{14pt}{14pt}\selectfont Track used: The End - theo5970}}
      \\[6pt]
      % Dòng cuối cùng: Thời gian diễn ra của tập (thường sẽ là ngày phát hành)
       & {\continuum\fontsize{16pt}{16pt}\selectfont Ngày phát hành: 07/07/2022}    \\[8pt]
    \end{tabular}
  \end{adjustwidth}
\end{table}
%=================================p==========NỘI DUNG==========================================
\color{black}

Ánh hoàng hôn nhuộm đỏ cả con sông nhỏ chảy qua thị trấn Aogami. Mặt nước lấp lánh, phản chiếu những cụm mây đang chuyển dần sang màu tím sẫm. Tiếng côn trùng bắt đầu cất lên từ cánh rừng phía xa, hòa với tiếng bước chân của nhóm học sinh đang đi dọc con đường đất nhỏ, ven bờ kè gạch rêu phong.

Năm cô gái - Hotaru, Mitsuki, Yuki, Kana và Kanade - đi phía trước, áo đồng phục đung đưa nhẹ trong gió, tiếng cười nói vang vọng khắp lối. Phía sau một quãng, Shino cùng ba người nhà Hikawa lặng lẽ bước theo, chỉ nghe chứ không xen vào.

Nàng tóc tím nghiêng đầu, khẽ chống cằm suy tư. ``Hmm\~{}\dots\ Hôm nay tôi thấy buổi thông báo trường nó cấn cấn thế nào ấy\dots''

Cô bạn thân quay sang, mái tóc trắng bạc hắt ánh chiều. ``Thế ư? Tôi thấy bà vẫn làm tốt như mọi khi thôi mà.''

Kana nhún vai, nhưng môi vẫn khẽ cong lên. ``Không biết nữa. Có lúc tôi thấy như mình đang đọc lại một bài diễn văn cũ thì đúng hơn.''

\img{e3-5a}

Hotaru khẽ đập vai Mitsuki, vừa đi vừa cười: ``Này nhé, biết vụ Mitsuki đi mua đồ hôm qua chưa? Bà ấy nha\dots''

Mitsuki trố mắt, gương mặt lập tức đỏ bừng. ``G-Gượm đã nào, Hotaru-san! Mình đã bảo cậu phải giữ bí mật chuyện đó mà\~{}\dots''

Yuki liền chen vào, giọng kéo dài như hát. ``Gì chứ? Chuyện đó nghe có vẻ hấp dẫn đấy\~{} Kể đi, kể đi!''

Cô nàng da sẫm đưa tay lên cằm, giả bộ suy nghĩ, môi nhếch thành nụ cười trêu chọc: ``Hmm\~{} Tôi có nên kể không ta\~{}''

Cô nàng tóc hồng gần như van nài: ``Đã bảo là đừng có kể mà\dots! Xin cậu đấy!''

\img{e3-5b}

Tiếng cười vang vọng, nhẹ và trong, hòa lẫn với tiếng gió lùa qua mặt sông. Cảnh tượng ấy thật bình dị, nhưng với Kuon, Kaigou và Suzuki, lại ẩn chứa một thứ cảm giác lạ lùng. Họ nhìn nhóm năm cô gái phía trước, giữa khung cảnh hoàng hôn ấy, và trong thoáng chốc, mọi thứ trở nên quen thuộc đến rợn người.

Kaigou khẽ lẩm bẩm: ``\vie{Y như lần đầu tiên\dots\ Họa chăng chỉ khác về nội dung cuộc nói chuyện thôi.}''

Kuon gật nhẹ, mắt nheo lại: ``\vie{Phải, ngay cả tiếng cười của họ\dots\ cũng không khác gì cả. Về cơ bản là vậy.}''

Suzuki không nói, chỉ siết chặt quai cặp trên vai, cảm giác như đang nhìn thấy một cảnh phim chiếu lại tới mức phát chán.

Bên cạnh họ, Shino đi lặng lẽ. Cô không mỉm cười, mă khuôn mặt có phần đanh lại. Ánh nhìn của cô thỉnh thoảng hướng về phía bên trái, nơi cánh rừng xanh thẫm đang chìm dần trong ánh chiều đỏ rực.

Bên trong cánh rừng ấy, cô biết có một ngôi đền cổ. Nơi mà sinh thời, cô đã từng cầu nguyện, rằng mình có thể kết được nhiều bạn, rằng cuộc sống mới này sẽ đầy ắp niềm vui.

Kuon, Kaigou và Suzuki nhìn theo ánh mắt ấy, lòng họ chợt nặng trĩu.

Họ biết nơi đó. Trong ký ức mờ nhạt của Kyuen, Saikyou và Reiko, đó là nơi mà Sora đã cầu khấn vị thần sau trận sạt lở.

Đó cũng là điểm cuối cùng của chuyến du ngoạn về quá khứ, nơi mà thứ ánh sáng màu trắng và tiếng chuông của ngôi đền vang lên.

Một cơn gió lạnh thổi qua, mang theo mùi đất ẩm và lá mục. Mặt sông phản chiếu ánh sáng nhấp nháy như đang vỡ vụn.

Cậu con trai khẽ thở ra, nhìn về phía trước, nơi Shino vẫn đang lặng lẽ dõi theo cánh rừng xa ấy, nét mặt cô trông vừa buồn vừa như đang cố hiểu điều gì đó không thể nói thành lời.

Suzuki nhìn cậu, rồi lắc đầu khẽ. ``\eng{Lần này\dots\ hy vọng có lẽ chính cô ấy sẽ là người nhận ra sớm.}''

``\eng{Đồng ý,}'' Game cất giọng. ``\eng{Nếu mọi chuyện cứ tiếp tục như vậy, chẳng biết điều gì sẽ xảy ra đâu.}''

Con đường dốc thoai thoải dẫn về khu dân cư. Những tiếng cười phía trước vẫn chưa dứt, vang vọng như bị kéo dài bất tận trong không khí đỏ quạch của hoàng hôn.

Mặt trời đã lặn hẳn phía sau dãy núi xa, chỉ còn lại sắc đỏ nhạt nhòa nhuộm loang trên mặt sông. Gió chiều trở nên lạnh hơn, mang theo hương ẩm của rêu và cỏ non ven bờ. Shino vẫn bước đi sau cùng, ánh mắt dõi theo từng bước chân mình trên con đường đất nhỏ, trong lòng dâng lên một cảm giác bất an mơ hồ.

Tất cả mọi thứ - tiếng gió, tiếng cười nói của nhóm bạn đi trước, thậm chí cả mùi hương thoảng qua trong không khí - đều quá hoàn hảo. Quá đều đặn. Như thể ai đó đã sắp đặt sẵn từng chi tiết nhỏ nhất để khung cảnh ấy không có lấy một vết nứt.

``\textit{Lạ thật\dots}'' cô thầm nghĩ, ``\textit{tại sao mình lại thấy như từng chuyện này\dots\ đã xảy ra rồi? Giống như một giấc mơ cũ đang lặp lại\dots}''

Tiếng Kana gọi khiến cô giật mình. ``Shino-san? Sao thế? Trông cậu như đang ngơ ngác ấy.''

Cô nàng học sinh mới chuyển ngẩng lên, nở một nụ cười gượng. ``Không có gì đâu. Chỉ là\dots\ tớ định ghé qua đền Aogami một lát trước khi về nhà.''

Kana chớp mắt, nhưng rồi nhún vai, giọng nhẹ tênh: ``Ờm, được thôi. Nhớ về sớm đó nha\~{}''

Cả nhóm năm cô gái tiếp tục đi về phía con dốc nhỏ, tiếng họ vang xa dần giữa ánh chiều tím. Khi những bóng lưng ấy đã khuất sau khúc quanh, Shino dừng lại trước bậc đá rêu phong dẫn lên ngôi đền cổ. Cô quay lại, nhìn ba người phía sau, và lần đầu tiên, giọng cô vang lên một cách khác lạ, lạnh và trầm hơn.

``Kuon-san, Kaigou-san, Suzuki-san. Đi cùng với tớ.''

Ba người thoáng sững lại. Trong ánh mắt Shino lúc này không còn vẻ ngây ngô của một cô gái mới chuyển trường nữa, mà là ánh nhìn tỉnh táo, lạnh lẽo như thể cô vừa tỉnh dậy sau một giấc mơ dài, giống như cách mà cô đã nói chuyện với họ trước đó.

Con đường dẫn lên đền thoai thoải, những bậc đá sứt mẻ lởm chởm, ẩn mình trong màn sương mờ buổi hoàng hôn. Ánh trăng vừa lên, đổ xuống lớp sáng bạc mong manh phủ trên vai họ. Mỗi bước chân vang lên khô khốc, vọng lại giữa không gian yên ắng, tựa như âm thanh của thời gian bị đóng băng.

Một cánh cổng torii đỏ sẫm hiện ra, uốn cong dưới vầng trăng tròn đang treo cao. Hai bên đường, những bụi hoa bỉ ngạn rũ xuống, đỏ rực như những đốm lửa tàn đang cháy trong bóng đêm. Cánh hoa khẽ rung rinh trong gió, vừa đẹp, vừa lạnh lẽo đến rợn người.

``\dots Lại là loài hoa này\dots'' giọng Suzuki cất lên, nhỏ như gió thoảng.

``Không thể nhầm được. Loài hoa nguyền rủa đó\dots\ vẫn còn đeo bám chúng ta đến tận nơi này,'' Kaigou chỉ nhìn lên cổng torii, đôi mắt đăm chiêu, im lặng bước tiếp.

Ngôi đền hiện ra trên đỉnh dốc. Bao quanh nó là rừng cây rậm rạp, thân cây cao lớn đan vào nhau thành những bóng đen sừng sững. Bậc đá dẫn vào đã nứt gãy, rêu xanh bám kín. Hai chiếc đèn đá cũ kỹ phủ bụi đứng chầu hai bên lối, ánh trăng chiếu lên khiến chúng trông như cặp mắt lặng lẽ dõi theo người đến.

Mái đền cong vút, chạm khắc những hoa văn cổ xưa đã mờ nhạt. Dưới chân họ, từng mảnh gỗ cũ kêu cọt kẹt yếu ớt, nhưng toàn bộ ngôi đền vẫn toát lên vẻ uy nghiêm lạnh lẽo, như một sinh thể đang ngủ yên giữa rừng sâu.

Shino dừng lại nơi bậc thềm, lặng nhìn lên ngôi đền ấy. Không có gì thay đổi kể từ sau bảy thập niên.

Từng chi tiết - từ màu ngói, dáng mái, đến vị trí hai chiếc đèn đá - đều giống hệt nơi cô từng đến cầu nguyện. Có thể có rêu, có cây mọc um tùm xung quanh theo thời gian, nhưng họ biết chắc chắn đó chính là ngôi đền đó.

Gió khẽ thổi qua, lay động tà váy đồng phục.

Trong khoảnh khắc ấy, cả ba người nhà Hikawa đều hiểu: điều mà họ đang thấy trước mắt, không phải là sự trùng hợp.

Một cơn gió nhẹ thoảng qua, mang theo tiếng xào xạc của lá khô. Một cơn gió lạnh, nhưng không phải cái lạnh của đêm, mà là sự rợn người, như thể thời gian bị ngưng đọng giữa hai thế giới.

\img{e3-5c}

Kaigou ngẩng lên, nhìn qua mái ngói phủ rêu xanh. Giọng cô khẽ, như nói với chính mình: ``\dots Không thay đổi gì cả.''

Suzuki cũng gật đầu, mắt dõi vào bậc đá nứt nẻ. ``Mặc cho thời gian đã trôi qua lâu như thế\dots\ Em có cảm giác rợn đến gai người, Onee-sama\dots''

Kuon im lặng. Cậu cảm nhận được điều gì đó đang dao động trong không khí, thứ năng lượng rất giống với ``nhiễu'' mà họ từng thấy trong các vòng lặp cũ.

Còn Shino\dots\ cô đứng yên trước cổng đền, đôi mắt ánh lên vẻ xa xăm.

``Tất cả những gì mọi người thấy\dots'' cô khẽ nói, ``\dots đều là giả. Mọi thứ. Cảnh vật, con người, tiếng cười\dots\ tất cả chỉ là những ảo ảnh mơ hồ mà thôi.''

Suzuki khẽ chau mày. ``Cậu biết rồi à?''

Shino gật đầu, nụ cười thoáng buồn. ``Tớ biết chứ. Tớ biết rõ chuyện gì đã xảy ra vào ngày hôm đó. Trận sạt lở định mệnh đó\dots\ đã cuốn tớ đi cùng ngôi trường, cùng cả những ký ức.''

Cô nắm chặt bàn tay, giọng lạc đi. ``Nhưng tớ vẫn muốn được sống thêm một chút nữa. Dù chỉ là trong giấc mơ này. Tớ muốn tiếp tục nói chuyện, cười đùa, đi học\dots\ như thể mọi thứ vẫn bình thường.''

Kaigou cúi đầu, giọng trầm: ``Đó không phải là sai, Shino-san. Ai cũng muốn được sống tiếp nếu còn cơ hội. Tôi hiểu mong ước của cậu mà. Nhưng\dots''

Kuon nhìn cô, ánh mắt hiền mà nghiêm, tiếp lời: ``Không sớm thì muộn, cậu cũng phải chấp nhận điều đó. Không ai có thể mãi trốn khỏi hiện thực, kể cả trong thế giới này.''

Game, từ chiếc đồng hồ trên cổ tay Kuon, cất giọng trầm tĩnh như một lời nhắc nhở: ``Ký ức chỉ là những mảnh vụn. Con người gắn bó với chúng vì sợ mất đi ý nghĩa của bản thân. Nhưng nếu mãi giữ chặt ảo ảnh, cô sẽ bị nó nuốt chửng.''

Shino mím môi. Cô hiểu rõ từng lời nói đó, nhưng ánh mắt vẫn không rời khỏi khoảng không phía sau cánh cổng đền. ``Tớ biết chứ\dots\ nhưng tớ chưa thể. Không phải bây giờ.''

Không ai nói thêm gì nữa. Chỉ còn lại tiếng gió lùa qua cánh rừng, và tiếng côn trùng kêu rền rĩ trong đêm.

Khi họ vẫn còn đang lặng im trong những suy nghĩ chồng chéo, một âm thanh khe khẽ vang lên. Ban đầu chỉ như tiếng gió, nhưng rồi dần rõ hơn, một tiếng nức nở, yếu ớt, thổn thức như của một đứa trẻ đang khóc trong tuyệt vọng.

``\dots Mình\dots\ Mình đã không thể làm được gì cả\dots''

Shino khẽ giật mình. ``Giọng nói đó\dots\ chẳng phải là\dots''

Suzuki tiếp lời, giọng căng thẳng: ``\dots Shinomiya-san?''

Âm thanh ấy dội lại giữa không gian tĩnh mịch, rồi đột nhiên, không khí quanh họ biến dạng. Cảnh vật méo mó như một tấm gương bị bẻ cong, từng mảng sương mờ trào dâng từ mặt đất. Làn sương dày đặc, lạnh lẽo, cuộn lấy từng bậc đá, từng mái ngói, từng tán cây, như thể muốn nuốt chửng cả ngôi đền cổ.

Khi làn sương tan dần, họ thấy một cô gái ngồi gục bên bậc đá, mặc bộ đồng phục cũ kỹ từ thế kỷ trước. Mái tóc hồng của cô phủ xuống, vai run rẩy theo từng tiếng nấc. Bộ đồng phục đen quen thuộc, nhưng đường may đã bạc màu, tay áo rách tơi tả, như thể vừa trải qua một cuộc chiến không tên. Đôi giày đen bùn đất dính đầy bùn đất, dấu hiệu của một hành trình dài dằng dặc không điểm dừng.

\img{e3-5d}

Shino bước lên, giọng run run: ``Shinomiya-san\dots cậu đang làm gì ở đây vậy?''

Cô gái ấy không đáp, chỉ siết chặt hai tay, vùi mặt trong lòng bàn tay. Những ngón tay tái nhợt, run rẩy, như thể đang cố giữ lại một điều gì đó đã mất từ lâu. Kuon và Kaigou nhìn nhau - trong bầu không khí nặng trịch ấy, giọng nói thanh lịch của Shinomiya Sakura dường như vừa thân thuộc, vừa xa lạ, như một hồn ma bị mắc kẹt giữa hai thực tại. Ánh mắt họ dán chặt vào bóng hình cô đơn ấy, như thể chỉ cần chớp mắt một cái, cô gái sẽ biến mất vào hư vô.

Suzuki khẽ gọi, ``Sakura-san\dots\ hãy nhìn bọn mình đi. Có chuyện gì thế?''

Shino cúi xuống, chạm nhẹ vào vai cô gái. Bàn tay cô run run, như thể đang chạm vào một bóng ma. Khoảnh khắc da thịt chạm nhau, một luồng điện lạnh chạy dọc sống lưng cô, khiến cô giật mình rụt tay lại.

Ngay khoảnh khắc ấy, một ánh sáng chói lòa lóe lên, nuốt trọn cả ngôi đền, cả khu rừng, và cả bốn người. Thế giới quanh họ tan rã như những mảnh kính vỡ, hòa vào trong luồng sáng trắng vô tận. Ánh sáng ấy mạnh đến mức khiến họ không thể mở mắt, khiến mọi âm thanh trở nên ù ù trong tai, khiến mọi cảm giác trở nên tê dại. Họ như đang rơi vào một khoảng không vô định, nơi không có trên, không có dưới, không có thời gian, không có không gian. Chỉ còn lại ánh sáng trắng và cảm giác bất lực tràn ngập tâm trí.

\dots

\dots

\begin{center}
  \uline{???}
\end{center}

Ánh sáng trắng dịu dần, nhường chỗ cho một thứ ánh lam mờ ảo. Không còn ngôi đền, không còn cánh rừng, chỉ còn một căn phòng thí nghiệm tĩnh lặng. Không khí nơi đây vương mùi ẩm của đất trồng xen lẫn mùi kim loại lạnh buốt. Dưới ánh đèn huỳnh quang nhấp nháy, những chậu cây non và dãy ống nghiệm lặng lẽ xếp hàng trên bàn thí nghiệm bằng thép không gỉ.

Trên kệ, vài bình thủy tinh chứa mẫu vật lơ lửng trong chất lỏng đục. Một vài cuốn sổ dày, cũ kỹ bị phủ bụi, nằm kề bên những thiết bị điện tử mà không ai trong họ từng thấy qua. Tiếng máy thông gió yếu ớt rền rĩ, như tiếng thở của một nơi đã bị lãng quên từ rất lâu.

\img{e3-5e}

Shino nhìn quanh, vẻ mặt đầy hoang mang. ``Đây là đâu đây\dots? Tại sao nơi này lại có nhiều cây cối như thế? Và\dots\ mấy cái máy móc lạ hoắc này là sao?''

Giọng trầm vang lên từ chiếc đồng hồ trên cổ tay Kuon: ``Trông nó giống như một phòng thí nghiệm sinh học.''

Kuon khẽ cau mày. ``Nhưng mà\dots\ tại sao lại ở đây chứ?''

Kaigou đặt tay lên bàn thép, cảm nhận lớp bụi lạnh. ``Có khi nào\dots\ chúng ta đang ở trong ký ức của Sakura-san không?''

Không ai dám chắc với câu trả lời đó, nhưng họ cũng không phủ định đó là một khả năng.

Cô nàng mắt xanh bước lại gần một chậu cây nhỏ, lá xanh non mơn mởn nhưng thân cây lại có những đường vân đỏ sẫm kỳ lạ.

``Cây này trông như thể đang được lai tạo giữa hai loài,'' Shino khẽ nói, tay chạm nhẹ vào chiếc lá. Chỉ một cái chạm, thân cây đã rung lên, tiết ra chút nhựa màu hổ phách, sền sệt và có mùi hăng nồng.

Cô chị lớn thì nhấc một ống nghiệm lên, bên trong là chất lỏng màu tím nhạt, có những hạt sáng li ti bơi lơi.

``Cấu trúc của chất này trông không giống bất kỳ hợp chất hữu cơ nào tôi từng thấy,'' Kaigou lắc nhẹ ống nghiệm, đôi mắt phản chiếu ánh sáng tím yếu ớt.

Cô em gái của cô thì mở một cuốn sổ trên kệ, trang đầu tiên ghi: \textit{``Thí nghiệm A-17: Độc tố lá bỉ ngạn - liều chết người 0,3 ml.''} Cô giật mình đóng sổ lại, mặt tái mét.

Còn cậu con trai, đứng trước một thiết bị như máy tính, màn hình đen nhưng vẫn phát ra tiếng bip bip đều đặn. ``Game-san, ông có thể truy cập vào nó không?''

Giọng Game vang lên, xen chút nhiễu tĩnh điện: ``Đang thử\dots\ Ồ, có một file ghi âm, nghe đi.''

Tiếng máy rè rè, rồi giọng Sakura vang lên, nhỏ và có gì đó mệt mỏi: \textit{``Ngày thứ 93: Tôi vẫn chưa tìm ra cách trung hòa độc tố. Nếu tôi không khẩn trương, cả thị trấn sẽ lại chìm trong biển máu. Tôi không muốn thêm một nạn nhân nào khác nữa\dots''}

*BÍP* Tiếng máy tắt ngấm.

Cả bốn người đứng lặng. Chỉ còn tiếng đèn huỳnh quang chập chờn, sáng lên rồi tắt dần, như hơi thở gấp gáp của ai đó vô hình.

Shino ôm vai mình, giọng run run: ``Cô ấy\dots\ đã sống một mình nơi này bao lâu để cứu mọi người?''

Kuon khẽ đặt tay lên vai cô, lắc đầu: ``Tôi cũng chẳng thể biết. Thậm chí tôi cũng không biết cô ấy đã làm gì trong cái phòng này nữa.''

Game bổ sung, giọng hiếm khi nghiêm túc: ``Phía cuối phòng có một cái tủ khoá bằng mã sinh trắc học. Nó cần dấu vân tay của Sakura-san để mở.''

Kaigou nhìn xuống bàn thí nghiệm, nơi có một chiếc găng tay cũ vẫn còn dính vết máu khô. ``Chắc phải dùng cái này thôi.''

Cả bốn nhìn nhau, cùng gật đầu.

Cô nàng tóc nâu run rẩy nhặt chiếc găng tay lên, đeo vào tay phải. Cô bước tới cánh cửa tủ lạnh kim loại xỉn màu, đặt lòng bàn tay lên bảng nhận dạng.

*TÍT*

Ánh đèn đỏ nhấp nháy một hồi, rồi chuyển sang xanh. Cánh cửa mở ra chậm rãi, phát ra âm thanh kéo dài rền rĩ như tiếng thở dài của chính căn phòng.

Bên trong là một chiếc hộp gỗ nhỏ, được đặt ngay ngắn giữa lớp bụi dày. Kaigou thận trọng nhấc nó ra, đặt lên bàn thí nghiệm.

``Một chiếc hộp mà lại khóa cẩn thận thế này ư, Onee-sama?'' Suzuki thì thầm.

Kuon mở nắp hộp, và cả bốn cùng lặng người.

Bên trong, không phải là tài liệu, không phải lọ hóa chất, mà là một nhúm hạt giống, được sắp ngay ngắn trên nền đất khô, phủ một lớp bụi mịn như tro tàn.

Shino cúi xuống, nhìn thật kỹ. Những hạt nhỏ xíu ấy ánh lên sắc đỏ nhạt như màu của hoa bỉ ngạn, nhưng không hẳn. Dường như có thứ gì đó đang khẽ chuyển động bên trong chúng.

``Hạt giống\dots?'' cô khẽ hỏi, giọng run run.

Game im lặng vài giây, rồi nói chậm rãi: ``Tôi cũng không chắc. Nhưng tôi cảm nhận được năng lượng trong đó. Nó không phải thứ bình thường. Có thể\dots\ chính chúng là `lời giải' mà Sakura-san đã để lại.''

Kaigou nhíu mày: ``Lời giải?''

``Nghe có vẻ mơ hồ nhỉ\dots'' Kuon đáp, ánh mắt lạnh đi.

Cả bốn người và cả giọng Game cùng im lặng, chỉ có tiếng đèn nhấp nháy thi thoảng chập chờn. Rồi, ở góc phòng, một vầng sáng yếu hắt xuống bàn. Dưới đó là một cuốn sổ bọc da nâu sẫm, bìa hơi rạn, mép giấy ngả vàng.

\img{e3-5f}

Cô nàng tóc nâu bước tới, khẽ nói: ``Cô ấy\dots\ muốn nói gì với chúng ta sao?''

Cô em út chau mày. ``Nhưng mà\dots\ đọc nhật ký riêng của người khác thì có hơi\dots''

Game lập tức ngắt lời, lạnh lùng: ``Thế còn việc đèn tự dưng nhấp nháy chỉ để soi vào cuốn đó, cô không thấy đáng nghi à?''

Một thoáng im lặng. Shino hít sâu, rồi cẩn thận nhấc cuốn nhật ký lên. Cô lật trang đầu tiên, nét chữ nghiêng, cứng cáp, đã mờ đi đôi chút.

\txtbx{Tôi là Shinomiya Sakura.

  \vspace{6pt}

  Tôi vẫn không thể nào quên được vụ sạt lở ngày đó, dù chỉ là chốc lát.}

Giọng đọc của Shino vang khẽ, nhưng dường như hòa lẫn vào không khí, như thể chính Sakura đang nói qua miệng cô.

Kaigou lặng người. ``Giọng văn này\dots\ nghe thật tĩnh, nhưng như chứa cả một cái gì đó như đang hối hận vậy.''

Shino tiếp tục lật trang.

\txtbx{Tất cả bắt đầu vào ngày hôm đó. Đêm thử thách lòng gan dạ định mệnh vào tháng Bảy, năm 1946.

  \vspace{6pt}

  Sau ngày hôm đó, cuộc sống bình thường của chúng tôi đã hoàn toàn bị đảo lộn.

  Tất cả đều đã bị đánh mất chỉ trong khoảnh khắc đột ngột đó. Chỉ còn lại mình tôi sống sót.

  \vspace{6pt}

  Những bạn học khác\dots\ hoặc là đã tử nạn, hoặc là mất tích không một dấu vết, trong đó có cả hai anh em sinh đôi nhà Haragen đó.

  \vspace{6pt}

  Cả Reiko-san cũng đột ngột biến mất, bặt vô âm tín.}

Suzuki siết nhẹ tay, giọng khàn đi: ``Hai anh em đó\dots\ chẳng phải là Onii-sama và Onee-sama sao? Em\dots\ cũng ở trong đó nữa.''

Kuon gật đầu chậm rãi. ``Cô ấy vẫn còn nhớ về chúng ta\dots\ cho tới tận cuối cùng.''

Shino nhìn trang giấy, giọng nghẹn lại: ``Cảm giác như thể cô ấy chưa bao giờ thoát khỏi ngày hôm đó.''

Cả bốn người tiếp tục đọc những trang tiếp theo trong cuốn nhật ký đó, cảm nhận được từng dòng chữ như thể Sakura đang thì thầm bên tai, kể lại nỗi đau và nước mắt đã thấm vào từng trang giấy úa vàng.

\txtbx{Trận sạt lở đất được cho là do những trận mưa lâu ngày, làm cho kết cấu đất bị xói mòn, cộng với nạn chặt phá rừng ở khu rừng Aogami.

  \vspace{6pt}

  Rất nhiều người đã lên tiếng phản đối, cho rằng vì chuyện đó mà tất cả học sinh đều ra đi mãi mãi.

  \vspace{6pt}

  Nhưng tôi biết lý do tận gốc xoay quanh cánh rừng mang tới điều bất hạnh này, cho những cây cối đã bị mục ruỗng bởi lời nguyền đó.}

Căn phòng chìm vào im lặng. Tiếng Shino khẽ vang lên, gần như là một lời độc thoại: ``Lời nguyền\dots\ của rừng Aogami\dots''

Kuon, Kaigou và Suzuki nhìn nhau, tất cả họ đều nhớ lại những gì Sakura từng kể với Mari.

Rằng thị trấn Aogami ngày xưa hay bị sạt lở.

Rằng dân phải hiến tế trẻ sống để xoa dịu thần núi.

Rằng từ đó núi ngừng lở nhưng hồn trẻ vẫn lảng vảng mang hận.

Và rằng tuyệt nhiên chớ bước vào rừng kẻo nghe tiếng gọi rồi chẳng thể trở lại như cũ.

Shino nheo mắt, vừa lướt qua lướt lại từng dòng, vừa nhớ lại những gì cô đã nghe: ``Vậy ra\dots\ Mari-san và Nanase-san phát điên là vì thứ đó.''

Suzuki khẽ gật. ``Không chỉ họ. Cả thị trấn Aogami này\dots\ có lẽ đã bị lời nguyền ấy nuốt chửng từ lâu.''

Ánh đèn chập chờn lần nữa. Một tiếng gió lạnh lùa qua khe cửa, làm mấy trang nhật ký khẽ rung lên, như muốn nói rằng câu chuyện ấy vẫn chưa kết thúc.

Cả bốn người lặng lẽ nhìn quanh căn phòng thí nghiệm. Ánh đèn huỳnh quang nhấp nháy, in những cái bóng loang lổ trên tường bê tông xám xịt. Mùi ẩm mốc và kim loại lạnh xộc vào mũi, khiến cả bốn người không khỏi khó chịu.

Shino ôm chặt cuốn nhật ký, mắt trố ra với những gì mà Sakura bộc bạch trong những trang giấy đã úa vàng ấy: ``Vậy\dots\ trong suốt khoảng thời gian sau đó, Shinomiya-san đã tự giam mình trong căn phòng này sao?''

\img{e3-5g}

Kuon gạt bụi trên chiếc kính hiển vi cũ, mắt vẫn quan sát từng ngóc ngách. ``Dường như là vậy. Cô ấy đã nắm chắc được nguyên nhân, vậy tôi nghĩ là cô ấy đã tìm cách giải quyết chuyện này.''

Suzuki khẽ vuốt lá cây trong chậu thí nghiệm, ánh mắt buồn bã. ``Sakura-san\dots\ dường như vẫn chưa thể quên sự vụ đó nhỉ?''

Game bật cười mỉa mai từ chiếc đồng hồ: ``Giống như ai đó đang cầm cuốn nhật ký đấy nhỉ.''

Shino giật mình, mặt đỏ bừng. Cô cúi xuống, giọng cười gượng gạo: ``Đ-Đúng vậy ha\dots''

Kaigou thở dài, nhìn về phía cánh cửa sắt phía cuối phòng. ``Dù sao\dots\ chúng ta nên tìm hiểu xem cô ấy đã làm gì ở đây. Có thể sẽ có manh mối về cách thoát khỏi ký ức này.''

\txtbx{Nếu cứ tiếp tục tàn phá rừng, những cơn sạt lở đất vẫn sẽ còn tiếp diễn. Nhưng nếu nó dừng lại, sẽ chẳng mấy chốc mà ta sẽ lại có thêm một Nanase-san tiếp theo\dots

  \vspace{6pt}

  Vậy, ta nên giải quyết như thế nào về cánh rừng này đây? Câu hỏi này vẫn luôn trằn trọc trong đầu tôi mãi.}

Shino khẽ đọc lên từng dòng, giọng cô run run. Trang giấy đã ngả màu, mực loang ra vì hơi ẩm, nhưng vẫn có thể thấy rõ những vết bút bị ấn mạnh như thể người viết đang cố kìm nén cảm xúc.

Căn phòng nhỏ trong ngôi đền chìm trong im lặng. Chỉ còn lại tiếng gió luồn qua khe gỗ, và tiếng thở từ chiếc đồng hồ của Kuon đều đặn.

Kaigou là người phá vỡ bầu không khí đầu tiên. ``Câu này\dots\ nghe như một lời tự trách. Có vẻ người viết đã chứng kiến tất cả, nhưng không thể làm gì để ngăn lại.''

Suzuki gật đầu, giọng cậu trầm xuống: ``Nếu họ biết trước được hậu quả, vậy nghĩa là thảm kịch này\dots\ đã được lặp đi lặp lại nhiều lần rồi.''

Shino khẽ siết cuốn nhật ký trong tay, đôi mắt ánh lên một nỗi buồn khó tả. ``Và có lẽ\dots\ cô ấy cũng bất lực, giống như tớ khi đó\dots''

Kuon im lặng hồi lâu rồi mới lên tiếng, giọng anh khàn khàn vì suy nghĩ: ``Nhưng `Nanase-san tiếp theo'\dots\ nghĩa là đã có nhueefu người từng chịu chung số phận như cô ấy? Có khi nào Sakura đã thấy điều gì đó mà không ai khác thấy?''

Trong đồng hồ, giọng Game vang lên, đều đều nhưng lạnh lẽo như thể vọng về từ một nơi xa lắm: ``Nếu ta hiểu đúng, thì người viết nhật ký đã biết rằng việc Nanase lâm vào tình trạng đó chỉ là hệ quả\dots\ của một vòng lặp. Dù rừng bị phá hay được giữ nguyên, thì bi kịch vẫn diễn ra. Một `Nanase' mới sẽ luôn xuất hiện.''

Cả bốn người đều im lặng. Họ nhìn nhau, và một nỗi rờn rợn len lỏi qua ánh mắt.

Shino nuốt khan, ngón tay run run lật sang trang kế tiếp. Tiếng giấy lật vang lên khẽ khàng trong không gian tĩnh lặng.

``\dots!''

Nhưng ngay khi trang kế tiếp mở ra, họ không còn thấy chữ, mà chỉ thấy một vệt đỏ nhòe, như máu đã thấm từ bên trong tờ giấy ra ngoài.

Ánh sáng nơi bàn thờ bỗng lập lòe. Một cơn gió mạnh thổi qua, cuốn tung những tờ giấy rơi tản mát khắp phòng.

Cuốn nhật ký khẽ rơi khỏi tay Shino, dừng lại ngay trước bậc thềm, mở ra một trang mới, nơi góc trang, một dòng chữ mờ hiện ra trong ánh sáng yếu ớt:

\txtbx{Tôi đã quyết định rồi.}

\end{document}